\section{Experimental Setup}
\label{sec:setup}

\subsection{Synthetic planted-alpha benchmark}
\label{sec:synthetic}

\paragraph{Purpose.} The benchmark checks that the harness works end to end
(evidence level~2 of the suite's governance ladder). It measures whether a
search method can recover signals that are \emph{known} to exist, and how
often the protocol selects anything on \emph{null} markets where nothing is
planted, under the same budget and protocol as the real-data study. It says
nothing about real markets.

\paragraph{Data-generating process.} For asset $i$ on day $t$, the return is
\[
r_{t,i}=\beta_i m_t+\sigma_i\varepsilon_{t,i}
      +s\sum_{k=1}^{K}w_k\,z_{k;t-\ell-1,i},\qquad
z_{k;t}=\mathrm{cs\_zscore}\big(g_k(\text{bars up to }t)\big),
\]
where:
\begin{itemize}
  \item each $g_k$ is a planted expression in the factor language;
  \item $s=\mathrm{SNR}\cdot\overline{\sigma}$, with $\overline{\sigma}$ the
    average idiosyncratic volatility;
  \item the weights satisfy $\sum_k w_k^2=1$.
\end{itemize}
The planted component of $r_t$ uses signal values from $t-\ell-1$. A signal
observed at close $t$ therefore earns its return exactly after the execution
lag used by the evaluator.

The remaining parameters are:
\begin{itemize}
  \item $\beta_i\sim U(0.5,1.5)$, $\sigma_i\sim U(0.012,0.030)$,
    $m_t\sim N(0.0003,0.010^2)$ and $\varepsilon_{t,i}\sim N(0,1)$;
  \item returns are clipped at $\pm0.5$;
  \item log volume follows an AR(1) process with persistence $0.7$ and
    innovation scale $0.35$, of which a share of $0.3$ is common across
    assets. Volume is independent of returns;
  \item open, high and low are noise around the close path.
\end{itemize}
The planted signals depend on prices that themselves contain the planted
component. The price path is therefore solved as a fixed point: returns are
rebuilt and every $z_k$ is re-evaluated \emph{with the package's own
evaluator} until the maximum absolute change is at most $10^{-12}$. On the
planted markets this took
\SynIterationsLow{} iterations at the low SNR level and \SynIterationsHigh{}
at the high one, with a maximum residual of \SynMaxResidual{}. The bars have finality \code{confirmed} and pass the
same panel validation as real data.

\paragraph{Planted signals.} There are $K=4$ signals with equal weights, in
three families:
\begin{itemize}
  \item \emph{textbook}: \code{reversal\_5}, \code{-ts\_mean(returns, 5)},
    and \code{abnormal\_volume\_20}, \code{volume / ts\_mean(volume, 20)};
  \item \emph{library family}: \code{volume\_return\_corr\_10},
    \code{-ts\_corr(returns, delta(log(volume), 1), 10)}, the volume--return
    correlation idea behind Alpha\#2 of \citet{kakushadze2016alphas}, whose
    formula is in the reference library;
  \item \emph{drawn}: \code{\SynDrawnName}, \SynDrawnExpression{}, drawn at
    random from the typed grammar by a procedure fixed before the draw. The
    first candidate meeting four criteria was accepted: depth at least 3,
    complexity at most 8 and a time-series operator; well defined on a
    reference null market; mean rank correlation below $0.3$ in absolute
    value with every library entry and every other planted signal; and a
    convergent fixed point at both SNR levels. The draw record, with every
    rejected candidate and its reasons, is committed with the results.
\end{itemize}
The textbook signals favour any proposer with prior knowledge of classic
factors, such as an LLM, and so does the library-family signal to a lesser
degree. Only the drawn signal is unknown to every proposer in advance; it
comes from the random baseline's own sampling distribution, which if anything
favours that baseline. Being drawn at random does not prove that the
expression is absent from the literature; the correlation criterion only
shows that it is far, in values, from the reference library. Measured on the
formation window of the planted markets, behavioural novelty against the
library was \SynNoveltyReversal{} and \SynNoveltyVolume{} for the textbook
signals, \SynNoveltyLibraryFamily{} for the library-family signal and
\SynNoveltyDrawn{} for the drawn one. Every planted signal (up to its sign,
which the protocol chooses) has a non-zero sampling probability under the
random-grammar baseline, and its components are on genetic programming's
menus.

\paragraph{Grid and protocol.} The grid and protocol settings are:
\begin{itemize}
  \item \SynSymbols{} fictional symbols and \SynSessions{} sessions;
  \item SNR $\in\{\SynSnrLow,\SynSnrHigh\}$ and market seeds $\{0,1,2\}$,
    plus \SynNullSeeds{} null markets (SNR $0$, seeds $0$ to $9$); the
    proposer seed is the market seed plus $10{,}000$;
  \item a budget of \SynBudget{} trials and batches of \SynBatch{};
  \item $\ell=h=1$, fractions $0.6/0.2/0.2$ and an embargo of 2 sessions;
  \item a formation screen (BH at $q=\SynFdrAlpha$ over behavioural
    classes, $m$ = budget minus behavioural duplicates), BH confirmation at
    $q=\SynConfirmAlpha$ on validation, top-$k=\SynTopK$ and decorrelation at
    $|\rho|<\SynMaxCorr$;
  \item exactly one reveal per run.
\end{itemize}
The realized dispersion ratio of the planted component to idiosyncratic
noise was \SynRealizedSnrLow{} on average at the low SNR level and
\SynRealizedSnrHigh{} at the high one.

\paragraph{Benchmark metrics.}
\begin{itemize}
  \item \emph{Recovery}: planted signal $k$ counts as recovered if some
    selected factor's oriented test-window values have mean cross-sectional
    Spearman correlation $\bar\rho\ge\SynRecoveryThreshold$ with $z_k$. The
    rule is signed: an oriented factor that bets against $z_k$ recovers
    nothing.
  \item \emph{Recall}: $|\bar\rho|\ge\SynRecoveryThreshold$ for any
    evaluated trial on the formation window, using every
    \SynRecallStride{}th formation date. Recall asks whether the proposer
    ever found the signal, selected or not.
  \item \emph{True FDP}: the share of selected factors that are false by
    ground truth, i.e.\ whose oriented rank correlation with the planted
    composite (proportional to the true expected return) on the test window
    is below \SynTrueDiscoveryThreshold{}. On a null market every selected
    factor is false.
  \item \emph{Test non-significance rate}: the share of selected factors
    whose one-sided test IC is not significant under BH at \SynTestAlpha{}
    within the selected set. It measures the power of the test window, not
    falsity.
  \item Composite and mean per-factor test IC and ICIR, behavioural and
    structural novelty, the selection funnel (behavioural duplicates,
    degenerate trials, screen survivors, confirmed candidates) and, on null
    markets, the share of runs that selected anything.
\end{itemize}
As an oracle reference, we score the planted signals themselves and their
composite on the same test window; the oracle's \emph{power} is the share of
planted signals that pass BH at \SynTestAlpha{} on the test window.

\subsection{Real-market study (pre-registered; not yet run)}
\label{sec:real}

The real-market design is fixed in the project's research plan before any
test-window access. We summarize it here.

\paragraph{Data.} US equity daily bars from MarketData, read only through the
suite's gateway with finality \code{confirmed}.
\todo{Document the provider's split and dividend adjustment convention and
the data-quality report (coverage, stale prices, extreme returns).}

\paragraph{Universes.}
\begin{itemize}
  \item U1: liquid large capitalizations, selected by median formation-window
    dollar volume.
  \item U2: a broader liquid universe.
\end{itemize}
Membership is frozen at the end of the formation window.
\todo{Final sizes and the share of names with truncated histories (a
survivorship diagnostic).}

\paragraph{Windows.} The data are split chronologically into formation,
validation and sealed test windows with the embargo rule of
Section~\ref{sec:metrics}. The test window must extend beyond the documented
training-data cutoff of every model evaluated. A prospective window of
sessions after the commitment date follows.
\todo{Exact dates, fixed before the first reveal.}

\paragraph{Arms and budgets.} The arms are:
\begin{itemize}
  \item A: random grammar;
  \item B: genetic programming;
  \item C: the classic library passed through the same selection;
  \item D: LLM-guided search;
  \item optionally E: an interactive-mining prompting
    style~\citep{wang2023alphagpt} re-implemented in the same language and
    budget;
  \item F: compute-matched A and B, whose budget matches the LLM arm's
    measured cost (at least $10^4$ trials) under the same trial accounting.
\end{itemize}
The primary budget is $B=200$ trials, with $R=5$ replicates per arm and
universe. Equal trial budgets are not equal compute, and GP and
reinforcement-learning alpha miners~\citep{yu2023alphagen} typically use far
more evaluations, which is why arm F exists. Every planned test reveal is
listed in advance and counted in the study registry.

\paragraph{Hypotheses and decision rules.}
\begin{itemize}
  \item \textbf{H1.} The LLM beats both baselines on the sealed test. Two
    estimands are tested separately: the mean daily difference in pooled
    composite test IC (market sampling; a stationary block bootstrap over
    test dates, in the spirit of \citet{hansen2005spa}) and the difference in
    replicate-level composite ICIR (search-procedure variability; an
    unpaired exact Mann--Whitney test, since LLM and baseline replicates share
    nothing to pair on). Holm correction over the two comparisons; support
    also requires a pre-set minimum median improvement.
  \item \textbf{H2 (synthetic, seeds 100--119).} The LLM recovers the drawn
    signal more often. The test is an exact McNemar test paired by market
    seed, with Holm correction. The seeds exclude the harness seeds already
    seen.
  \item \textbf{H3.} The LLM's edge is not explained by period memorization.
    The decision uses only the difference-in-differences $D$ of pre- versus
    post-cutoff test IC between the LLM and the pre-specified GP baseline:
    memorization is flagged if a bootstrap interval for $D$ lies below zero,
    and absence of memorization is declared only by an equivalence test with
    a margin fixed in advance. It is declared underpowered with fewer than
    250 post-cutoff sessions.
  \item \textbf{H4.} Formation feedback improves search. It is tested on
    synthetic recall and on validation-window ICIR only.
\end{itemize}
Every test is one-sided at level $0.05$; the equivalence test of H3 is a pair of one-sided tests.

\paragraph{Ablations.} Ablations use validation data only. They vary:
\begin{itemize}
  \item feedback (on or off);
  \item the rationale requirement (on or off);
  \item anonymization (on or off, as a contamination probe in a pre-cutoff
    window only);
  \item reasoning effort;
  \item budget.
\end{itemize}

\subsection{Implementation and reproducibility}

\paragraph{Software.} Everything is implemented in Python. The recorded
software stack is \SynSoftware{}. The offline test suite checks:
\begin{itemize}
  \item absence of look-ahead for every operator;
  \item the information barrier;
  \item ledger tamper detection;
  \item the seal rules;
  \item the LLM backend's request shape and error handling, with a fake
    client.
\end{itemize}

\paragraph{Runs and artifacts.} The synthetic benchmark (\SynRuns{} planted
and \SynNullRuns{} null runs) took \SynRuntime{}~s of wall-clock time. It is
deterministic: ledger timestamps are disabled, and rerunning four of its runs
under a different hash seed reproduced their ledger heads and every scored
field. Each run stores:
\begin{itemize}
  \item its configuration and the hash-chained ledger;
  \item the selection and the reveal;
  \item a provenance file with the data hash, template hashes, served model
    identifiers, software versions and the ledger head.
\end{itemize}
LLM runs additionally record every request and its outcome, failures
included, for exact offline replay.
